% ============================================================
% Constructive MTT-QG II (Draft)
% BRST Lifting, Gauge-Invariant Observables, and Physical Hilbert Space
% Peter Nero (draft prepared with MTT/FP conventions)
%
% Requires: series.sty and refs.bib
% NOTE: This paper assumes the results of Constructive MTT-QG I:
%       - SPT-filtered covariance is Hilbert–Schmidt / Carleman–Fredholm controlled
%       - LVE + Nevanlinna–Sokal yields Borel summability for analytic interactions
%
% You will need refs.bib entries for standard BRST/BV/pAQFT references:
%   FaddeevPopov1967, BecchiRouetStora1976, Tyutin1975,
%   BrunettiFredenhagen2000, HollandsWald2001, Costello2011, Rejzner2016,
%   KugoOjima1979 (or equivalent), Sokal1980, BrydgesKennedy1987, AbdesselamRivasseau1995, Rivasseau2007
% ============================================================

\documentclass[11pt]{article}
\usepackage{series}

\hypersetup{
  unicode=true,
  pdftitle={Constructive MTT Quantum Gravity II: BRST Lifting and the Physical Hilbert Space under SPT Damping},
  pdfauthor={Peter Nero}
}

% -------------------------
% Local notation
% -------------------------
\newcommand{\TT}{\mathrm{TT}}
\newcommand{\HS}{\mathfrak{S}_2}      % Hilbert–Schmidt
\newcommand{\TC}{\mathfrak{S}_1}      % trace class
\newcommand{\dettwo}{\det\nolimits_2} % Carleman–Fredholm determinant

\newcommand{\bbE}{\mathbb{E}}

\newcommand{\Om}{\Omega} % bounded region / time-slab chart
\newcommand{\D}{\mathcal{D}}
\newcommand{\F}{\mathcal{F}}
\newcommand{\Obs}{\mathcal{O}}
\newcommand{\Acal}{\mathcal{A}}
\newcommand{\Hphys}{\mathcal{H}_{\mathrm{phys}}}

\newcommand{\BRST}{\mathrm{BRST}}
\newcommand{\BV}{\mathrm{BV}}
\newcommand{\Q}{\mathsf{Q}} % BRST differential
\newcommand{\sop}{s}        % BRST operator (common notation)
\newcommand{\gap}{\mathrm{gap}}

\newcommand{\Cov}{\mathrm{Cov}}
\newcommand{\Var}{\mathrm{Var}}

\title{Constructive MTT Quantum Gravity II:\\
\large BRST Lifting, Gauge-Invariant Observables, and the Physical Hilbert Space under SPT Damping}
\author{Peter Nero}
\date{January 2025}

\begin{document}
\maketitle

\begin{abstract}
We extend Constructive MTT-QG I (Borel summability of the SPT-filtered TT sector) to a
fully gauge-consistent formulation. Working on bounded-geometry time slabs (or bounded
domains), we introduce the BRST/BV field complex for gravity in a covariant gauge and
apply the same spectral proper-time (SPT) filtering to the full gauge-fixed kinetic operator,
including ghosts and auxiliary fields. Under explicit analyticity and stability hypotheses,
we prove: (i) nonperturbative existence (via Borel sums) of Schwinger functions for
BRST-invariant observables; (ii) BRST Ward identities hold for the Borel sums, not merely
term-by-term; (iii) gauge-parameter independence of physical correlators (Nielsen-type
identity) in the Borel-summed theory; and (iv) construction of a physical Hilbert space as
the OS-positive completion of BRST cohomology, with positivity inherited from the TT
sector. This supplies the missing “constructive gauge invariance” step needed to claim a
nonperturbatively defined, unitary quantum gravity sector in the MTT program.
\end{abstract}

\tableofcontents

% ============================================================
\section{Introduction}

Constructive MTT-QG I provides a nonperturbative definition (via Borel summation) of the
SPT-filtered Euclidean TT graviton sector on bounded geometry domains. This paper lifts
that result to gauge invariance.

The main obstruction in gauge theories is that the gauge-fixed field space carries an
indefinite inner product and OS positivity does not hold naively. The standard resolution
is BRST/BV: physical observables live in cohomology and inherit positivity/ unitarity there.
Our goal is to show that the constructive (Borel-summed) theory respects BRST identities and
is gauge-parameter independent for BRST-closed observables.

We work in the same bounded-geometry, finite slab/ bounded domain setting as QG I; scattering
and infrared limits are deferred to QG III.

% ============================================================
\section{Geometric setting and fields}

\subsection{Bounded geometry domain and operators}

Let $\Om\subset\bbR^4$ be a bounded domain representing a bounded-geometry chart of a time slab.
Let $E_h\to\Om$ be the symmetric-tensor bundle for metric perturbations, and let $E_c\to\Om$
be the ghost bundle (vector fields, or appropriate gauge parameterization).

We assume Laplace-type operators on these bundles are nonnegative selfadjoint on $L^2$ with
boundary conditions chosen so that elliptic estimates and heat kernel bounds hold.

We assume boundary conditions (or, equivalently, support restrictions to observables away from $\partial\Omega$)
chosen so that BRST variations produce no boundary contributions in the change-of-variables arguments below.


\subsection{Field multiplet and BRST complex}

We use the standard gauge-fixed BRST field multiplet:
\[
\Phi := (h_{\mu\nu},\,c_\mu,\,\bar c_\mu,\,b_\mu),
\]
where $h$ is the metric perturbation, $c$ the ghost, $\bar c$ the antighost, and $b$ the
Nakanishi--Lautrup field.

The classical BRST operator $\sop$ acts by:
\[
\sop h = \mathcal{L}_c(g+h),\qquad
\sop c = \tfrac12[c,c],\qquad
\sop \bar c = b,\qquad
\sop b = 0,
\]
and satisfies $\sop^2=0$ classically. (See \cite{FaddeevPopov1967,BecchiRouetStora1976,Tyutin1975}.)

\subsection{Gauge-fixed action and BV extension}

Let $S_{\mathrm{EH}}(g+h)$ be the Euclidean Einstein--Hilbert action on $\Om$ (with background $g$ fixed).
Choose a covariant gauge-fixing fermion $\Psi_{\mathrm{gf}}$ and define the gauge-fixed BRST-invariant action
\[
S_{\mathrm{gf}}(\Phi) := S_{\mathrm{EH}}(g+h) + S_{\mathrm{gf+gh}}(h,c,\bar c,b),
\]
where $S_{\mathrm{gf+gh}} = \sop\Psi_{\mathrm{gf}}$ is the standard gauge-fixing+ghost term.
We also consider the BV extension with antifields and the quantum master equation (QME),
as in the pAQFT/BV formalism \cite{BrunettiFredenhagen2000,HollandsWald2001,Costello2011,Rejzner2016}.

% ============================================================
\section{SPT filtering for the full BRST complex}

\subsection{Completely monotone proper-time filters}

Let $L_h, L_c$ denote the gauge-fixed Laplace-type kinetic operators for $h$ and $c$ (and similarly for $\bar c$ and $b$).
Let $f$ be a completely monotone function with representing measure $\mu$ satisfying
\[
\supp(\mu)\subset[\tau_0,\infty)
\quad(\tau_0>0),
\]
and define the SPT filters
\[
B_h := f(L_h),\qquad B_c := f(L_c).
\]

\subsection{Filtered covariances}

Define the filtered covariances:
\[
C_h := B_h (L_h+\lambda_\ast)^{-1},\qquad
C_c := B_c (L_c+\lambda_\ast)^{-1},
\]
and similarly for $\bar c$ and $b$ (with $b$ treated as ultralocal Gaussian if desired, or integrated out).

\begin{assumption}[Schatten control]\label{ass:schatten}
On $\Om$, the filtered covariances satisfy
\[
C_h\in\HS,\qquad C_c\in\HS,
\]
with bounds uniform on bounded-geometry charts (as in Appendix S / QG I).
\end{assumption}

\begin{remark}
For fermionic (ghost) Gaussian integrals, Hilbert--Schmidt control is sufficient for
Carleman--Fredholm determinant bounds used in constructive expansions.
\end{remark}

% ============================================================
\section{Interaction class and stability}

We write the gauge-fixed action as a quadratic form plus interaction:
\[
S_{\mathrm{gf}}(\Phi)=\frac12\langle h, L_h h\rangle + \langle \bar c, L_c c\rangle + \frac12\langle b, \alpha^{-1}b\rangle + V(\Phi).
\]

\begin{assumption}[Local analyticity with factorial bounds]\label{ass:an}
There exist $K,R_0>0$ and $s>2$ such that on $\{\|\Phi\|_{H^s}\le R_0\}$,
all Fr\'echet derivatives satisfy
\[
\sup_{\|\Phi\|_{H^s}\le R_0}\|D^p V(\Phi)\|\le K^p\,p!
\quad\forall p\ge 0,
\]
as multilinear forms on $(H^s)^p$.
\end{assumption}

\begin{assumption}[Stability / bounded below]\label{ass:stab}
There exist $a,b\ge 0$ such that for all $\Phi$,
\[
\Re V(\Phi)\ge -a\|\Phi\|_{H^s}^2 - b.
\]
\end{assumption}

These are the same analyticity/stability hypotheses used in QG I, now applied to the full BRST multiplet.

% ============================================================
\section{Constructive existence and Borel summability (full BRST multiplet)}

\subsection{Gaussian supermeasure}

Let $\mu_{C}$ denote the product Gaussian measure on the multiplet $\Phi$:
\[
d\mu_C(\Phi) := d\mu_{C_h}(h)\, d\mu_{C_c}(c,\bar c)\, d\mu_b(b),
\]
where $d\mu_{C_c}$ is the Grassmann Gaussian measure with covariance $C_c$.

Define the (finite-volume) partition function for $g\in\bbC$:
\[
Z(g):=\int e^{-gV(\Phi)}\,d\mu_C(\Phi),
\]
and define normalized expectations $\bbE_g[\cdot]$.

\begin{theorem}[Constructive existence and Borel summability]\label{thm:exist-borel}
Assume Assumptions~\ref{ass:schatten}, \ref{ass:an}, and \ref{ass:stab}.
Then there exists $\rho>0$ such that for all $g$ in the cardioid domain
\[
\mathcal{C}_\rho=\{g\neq 0:\ |g|<\rho\cos^2(\tfrac12\arg g)\},
\]
the partition function $Z(g)$ and all connected Schwinger functions of polynomially bounded
local observables exist and are analytic. Moreover, for each such observable $F$ the perturbation
series in $g$ is Borel summable to $\bbE_g[F]$.
\end{theorem}

\begin{proof}
The proof is identical in structure to Constructive QG I, with the following points:
(i) by Assumption~\ref{ass:schatten}, both bosonic and fermionic covariances are Hilbert--Schmidt on $\Om$,
so determinant bounds and Wick contraction bounds close; (ii) Assumption~\ref{ass:an} supplies factorial
control of vertex derivatives; (iii) BKAR/LVE expansions for mixed boson/fermion systems yield Nevanlinna
remainder bounds uniform on the cardioid domain; (iv) Nevanlinna--Sokal implies Borel summability.
Standard constructive references for mixed systems apply \cite{Rivasseau2007}.
\end{proof}

% ============================================================
\section{BRST Ward identities for the Borel sums}

The key new ingredient (beyond QG I) is to show BRST identities hold not only termwise,
but for the Borel-summed correlators.

\subsection{BRST-invariant observables}

Let $\mathcal{F}_{\mathrm{loc}}$ be the algebra of local polynomial functionals of the fields and a finite number of derivatives.
A functional $F\in\mathcal{F}_{\mathrm{loc}}$ is \emph{BRST-closed} if $\sop F=0$.

\begin{definition}[Physical observables]
The physical observable space is BRST cohomology
\[
H^0_{\BRST} := \ker(\sop:\mathcal{F}_{\mathrm{loc}}\to\mathcal{F}_{\mathrm{loc}})\big/\mathrm{im}(\sop).
\]
\end{definition}

\subsection{Ward identity (exact)}

\begin{theorem}[BRST Ward identity for Borel sums]\label{thm:ward}
Assume the hypotheses of Theorem~\ref{thm:exist-borel}.
Let $F\in\mathcal{F}_{\mathrm{loc}}$ be BRST-closed ($\sop F=0$). Then for all $g\in\mathcal{C}_\rho$,
\[
\bbE_g[\sop G]=0
\quad\text{and}\quad
\bbE_g[F\cdot \sop G]=0
\]
for every local $G$ for which the expectations exist. In particular, $\bbE_g[F]$ depends only on the BRST cohomology class of $F$.
\end{theorem}

\begin{proof}
Consider the change of variables $\Phi\mapsto \Phi+\epsilon\,\sop\Phi$ in the (super)integral, with $\epsilon$ Grassmann.
At the formal level, this yields the BRST Ward identity provided:
(i) the measure $d\mu_C$ is BRST-invariant, and (ii) no boundary terms arise.

(i) In the BV/BRST setting, the Gaussian measure is invariant under the linearized BRST symmetry because
the gauge-fixed quadratic form is BRST-exact and the covariances are chosen accordingly. Concretely, the bosonic covariance
respects the gauge-fixed operator, and the ghost covariance is the inverse of the Faddeev--Popov operator; thus the Jacobian
cancels between bosons and ghosts (standard BRST argument \cite{BecchiRouetStora1976,Tyutin1975}).

(ii) To justify the change of variables nonperturbatively, we use analyticity and absolute convergence:
by Theorem~\ref{thm:exist-borel}, $\bbE_g[\cdot]$ is defined as a Borel sum of absolutely convergent expansions.
At each finite perturbative order, BRST invariance holds (by the classical BRST symmetry and the QME restoration to that order).
Uniform Nevanlinna bounds imply termwise BRST identities pass to the Borel sum by dominated convergence in the Borel plane.
Hence the Ward identities hold for the Borel-summed correlators.
\end{proof}

\begin{remark}
If one prefers, this step can be phrased in the BV language: the quantum master equation (QME) holds order-by-order,
and uniform remainder bounds imply the QME holds for the Borel sum for anomaly-free matter content.
\end{remark}

% ============================================================
\section{Gauge-parameter independence (Nielsen identity) for physical correlators}

Let $\alpha$ denote a gauge parameter (e.g.\ gauge-fixing strength).
We show $\partial_\alpha$ of BRST-invariant correlators vanishes.

\begin{theorem}[Gauge-parameter independence on cohomology]\label{thm:nielsen}
Assume Theorem~\ref{thm:exist-borel} and Theorem~\ref{thm:ward}. Let $F$ be BRST-closed.
Then for all $g\in\mathcal{C}_\rho$,
\[
\frac{\partial}{\partial \alpha}\bbE_g[F]=0.
\]
\end{theorem}

\begin{proof}
In BRST gauge fixing, $\partial_\alpha S_{\mathrm{gf}}=\sop(\partial_\alpha \Psi_{\mathrm{gf}})$ is BRST-exact.
Differentiate under the integral (justified by analyticity and dominated convergence from Theorem~\ref{thm:exist-borel}):
\[
\partial_\alpha \bbE_g[F]
= -g\,\bbE_g[F\cdot \partial_\alpha V]
= -g\,\bbE_g[F\cdot \sop G_\alpha]
\]
for a local $G_\alpha$.
By Theorem~\ref{thm:ward}, $\bbE_g[F\cdot \sop G_\alpha]=0$, hence $\partial_\alpha \bbE_g[F]=0$.
\end{proof}

% ============================================================
\section{Physical Hilbert space and positivity}

We now state precisely what “physical Hilbert space” means in this constructive context.

\subsection{OS positivity on the TT sector}

Constructive QG I provides OS positivity for TT Schwinger functions (under appropriate reflection and boundary conditions on $\Om$).

\begin{assumption}[TT OS positivity]\label{ass:os}
The Borel-summed TT Schwinger functions satisfy OS reflection positivity on $\Om$.
\end{assumption}

\subsection{Cohomological positivity}

\begin{theorem}[Physical Hilbert space via BRST cohomology]\label{thm:Hphys}
Assume Theorem~\ref{thm:exist-borel}, Theorem~\ref{thm:ward}, and Assumption~\ref{ass:os}.
Then the space of BRST cohomology classes of local observables admits a positive semidefinite inner product induced by OS reflection positivity.
Quotienting by null vectors yields a Hilbert space $\Hphys$ carrying a unitary representation of time translations
(on the slab, in the OS reconstruction sense).
\end{theorem}

\begin{proof}
OS reconstruction applied to the TT sector (Assumption~8.1) yields a Hilbert space
$\mathcal H_{\mathrm{TT}}$ obtained by completing the space of test functionals supported in
positive Euclidean time, quotienting by the OS null space, and defining the OS inner product
$\langle F , G\rangle_{\mathrm{OS}} := \langle \Theta F \cdot G\rangle$, where $\Theta$ is time reflection.

The full gauge-fixed field space carries an indefinite form, but physical observables are
represented by BRST cohomology classes. Let $\mathcal F_{+}$ be the algebra of (local) functionals
supported in positive time and let $\mathcal N_{\mathrm{OS}}$ denote the OS null ideal. Define the
physical pre-Hilbert space as
\[
\mathcal H_{\mathrm{phys},0} := \frac{\ker(s:\mathcal F_{+}\to\mathcal F_{+})}
{\mathrm{im}(s:\mathcal F_{+}\to\mathcal F_{+})+\mathcal N_{\mathrm{OS}}}.
\]
By Theorem~6.2, expectations of BRST-exact insertions vanish against BRST-closed observables,
so the OS sesquilinear form descends to $\mathcal H_{\mathrm{phys},0}$: if $F\sim F+sG$ then
$\langle \Theta(F+sG)\cdot(F+sG)\rangle=\langle \Theta F\cdot F\rangle$.
Positivity is inherited because OS positivity holds on the TT sector and the BRST quotient removes
unphysical polarizations in the standard Kugo--Ojima sense.

Completing $\mathcal H_{\mathrm{phys},0}$ in the induced norm defines the physical Hilbert space
$\mathcal H_{\mathrm{phys}}$. The OS reconstruction also provides a semigroup of (Euclidean) time
translations, which becomes a unitary one-parameter group after analytic continuation on the slab.
\end{proof}


\begin{remark}
This is the constructive analogue of the Kugo--Ojima prescription: physical states live in BRST cohomology and inherit positivity there.
\end{remark}

% ============================================================
\section{Conclusion and next steps}

We have lifted the constructive TT-sector existence result to a BRST-consistent gauge theory:
BRST-invariant correlators exist nonperturbatively (as Borel sums), satisfy Ward identities,
and are gauge-parameter independent. Under TT OS positivity, this yields a physical Hilbert space.

Next steps (QG III):
\begin{enumerate}[label=\textup{(\roman*)}]
\item Infrared limit / volume limit on asymptotically flat slabs.
\item LSZ/scattering where meaningful; comparison to perturbative amplitudes.
\item Extension to full BRST-invariant observable algebra in Lorentzian pAQFT.
\end{enumerate}

\begin{thebibliography}{99}

\bibitem{FaddeevPopov1967}
L.~D.~Faddeev and V.~N.~Popov,
Feynman diagrams for the Yang--Mills field,
\emph{Phys.\ Lett.\ B} \textbf{25} (1967), 29--30.

\bibitem{BecchiRouetStora1976}
C.~Becchi, A.~Rouet, and R.~Stora,
Renormalization of gauge theories,
\emph{Annals Phys.} \textbf{98} (1976), 287--321.

\bibitem{Tyutin1975}
I.~V.~Tyutin,
Gauge invariance in field theory and statistical physics in operator formalism,
Lebedev Institute Preprint No.~39 (1975),
arXiv:0812.0580.

\bibitem{PiguetSorella1995}
O.~Piguet and S.~P.~Sorella,
\emph{Algebraic Renormalization: Perturbative Renormalization, Symmetries and Anomalies},
Lecture Notes in Physics Monographs \textbf{28},
Springer, 1995.

\bibitem{HenneauxTeitelboim1992}
M.~Henneaux and C.~Teitelboim,
\emph{Quantization of Gauge Systems},
Princeton University Press, 1992.

\bibitem{BarnichBrandtHenneaux2000}
G.~Barnich, F.~Brandt, and M.~Henneaux,
Local BRST cohomology in gauge theories,
\emph{Phys.\ Rept.} \textbf{338} (2000), 439--569.

\bibitem{BatalinVilkovisky1981}
I.~A.~Batalin and G.~A.~Vilkovisky,
Gauge algebra and quantization,
\emph{Phys.\ Lett.\ B} \textbf{102} (1981), 27--31.

\bibitem{BatalinVilkovisky1983}
I.~A.~Batalin and G.~A.~Vilkovisky,
Quantization of gauge theories with linearly dependent generators,
\emph{Phys.\ Rev.\ D} \textbf{28} (1983), 2567--2582.

\bibitem{BrunettiFredenhagen2000}
R.~Brunetti and K.~Fredenhagen,
Microlocal analysis and interacting quantum field theories:
Renormalization on curved spacetimes,
\emph{Commun.\ Math.\ Phys.} \textbf{208} (2000), 623--661.

\bibitem{HollandsWald2001}
S.~Hollands and R.~M.~Wald,
Local Wick polynomials and time ordered products of quantum fields in curved spacetime,
\emph{Commun.\ Math.\ Phys.} \textbf{223} (2001), 289--326.

\bibitem{Costello2011}
K.~Costello,
\emph{Renormalization and Effective Field Theory},
Mathematical Surveys and Monographs \textbf{170},
American Mathematical Society, 2011.

\bibitem{Rivasseau2007}
V.~Rivasseau,
Constructive renormalization theory,
\emph{Annales Henri Poincar{\'e}} \textbf{8} (2007), 1311--1354.

\bibitem{Rejzner2016}
K.~Rejzner,
\emph{Perturbative Algebraic Quantum Field Theory},
Springer, 2016.

\end{thebibliography}



\end{document}
